\section{Discussion}\label{sec:discussion}

\subsection{The searched sequence saves 2,866.4 to 3,120.5 weighted minutes per period at \texorpdfstring{$\nu = 0.2$}{nu = 0.2}, at the cost of a hand-traceable rule}\label{sec:recommendation}

\sumpara We recommend the searched sequence and state its benefit, its cost, and the condition that would reverse the choice. \sumparaend
\fin We recommend that the planner sequence the line with iterated greedy. \fine
\fin For a period like the 45-job instance with processing times varying at a coefficient of variation of 0.2, its fixed plan lowers mean total weighted tardiness by 3,120.5 weighted minutes (standard error 150.1) on the main scenarios and by 2,866.4 (standard error 121.3) on an independent second set, leaves 11.5~jobs late per period on average against 36.1, and cuts mean waiting between stages by 48.0~percent. \fine
\fin The planner gives up a rule that can be traced by hand, accepts a result that depends on a seed and takes seconds to compute, 3.22 at 45~jobs and 37.2 at 120, and on very small batches may accept more waiting, as on the 6-job instance. \fine
\fin The choice needs a fresh comparison if processing times vary much more than at $\nu = 0.2$: at $\nu = 0.8$ the mean gain held at 3,329.0 weighted minutes, but the 10th percentile of the per-period gain fell to 851.5 and the rule's plan won 1 of 30~scenarios, so a planner who weighs the worst periods most should repeat the comparison at the line's measured variability. \fine

\subsection{The evidence covers one generator, independent noise, and a plan that is never revised}\label{sec:scope}

\sumpara Our design bounds the claim in four ways that a planner should check against the line. \sumparaend
\fin All four instances come from one generator, and the stress test covers only the 45-job instance; the instances have processing times of 2 to 20~minutes and median due-date windows of 154 to 158~minutes, so a line with tighter due dates or more uneven stages may show different margins. \fine
\fin Our scenarios draw each of the 360 processing times of the 45-job instance independently, so they do not represent a common cause, such as a slow shift or a degraded machine, that lengthens many operations at once. \fine
\fin The planner in our study never revises the sequence, which isolates the value of the plan itself and leaves out any gain from rescheduling when actual times are known. \fine
\fin Optimality is proven only up to 15~jobs, and on the 120-job instance other seeds found sequences up to 5.1~percent better, so the gains we report for the largest instance understate what the search can deliver. \fine

\subsection{Three extensions follow from these limits}\label{sec:next}

\sumpara Each limit points to one extension with a clear test. \sumparaend
\fin A search that scores each candidate sequence by its mean tardiness over sampled scenarios would target the variability directly and could narrow the relative gap at high $\nu$, which the sweep of \cref{sec:sweep} can measure. \fine
\fin Scenarios with correlated factors, in which one draw lengthens a whole stage or a whole shift, would test whether the absolute gain survives a common cause. \fine
\fin Re-running the search at each release with the times already realized would show how much of the stress-test loss a revised plan recovers. \fine
